\documentclass{amsart}
% For dealing with references we use the comment environment
\usepackage{verbatim}
\newenvironment{reference}{\comment}{\endcomment}
%\newenvironment{reference}{}{}
\newenvironment{slogan}{\comment}{\endcomment}
\newenvironment{history}{\comment}{\endcomment}

% For commutative diagrams we use Xy-pic
\usepackage[all]{xy}

% We use 2cell for 2-commutative diagrams.
\xyoption{2cell}
\UseAllTwocells

% We use multicol for the list of chapters between chapters
\usepackage{multicol}

% This is generally recommended for better output
\usepackage{lmodern}
\usepackage[T1]{fontenc}

% For cross-file-references

% Package for hypertext links:
\usepackage{hyperref}

% For any local file, say "hello.tex" you want to link to please

% Theorem environments.
%
\theoremstyle{plain}
\newtheorem{theorem}[subsection]{Theorem}
\newtheorem{proposition}[subsection]{Proposition}
\newtheorem{lemma}[subsection]{Lemma}

\theoremstyle{definition}
\newtheorem{definition}[subsection]{Definition}
\newtheorem{example}[subsection]{Example}
\newtheorem{exercise}[subsection]{Exercise}
\newtheorem{situation}[subsection]{Situation}

\theoremstyle{remark}
\newtheorem{remark}[subsection]{Remark}
\newtheorem{remarks}[subsection]{Remarks}

\numberwithin{equation}{subsection}

% Macros
%
\def\lim{\mathop{\mathrm{lim}}\nolimits}
\def\colim{\mathop{\mathrm{colim}}\nolimits}
\def\Spec{\mathop{\mathrm{Spec}}}
\def\Hom{\mathop{\mathrm{Hom}}\nolimits}
\def\Ext{\mathop{\mathrm{Ext}}\nolimits}
\def\SheafHom{\mathop{\mathcal{H}\!\mathit{om}}\nolimits}
\def\SheafExt{\mathop{\mathcal{E}\!\mathit{xt}}\nolimits}
\def\Sch{\mathit{Sch}}
\def\Mor{\mathop{\mathrm{Mor}}\nolimits}
\def\Ob{\mathop{\mathrm{Ob}}\nolimits}
\def\Sh{\mathop{\mathit{Sh}}\nolimits}
\def\NL{\mathop{N\!L}\nolimits}
\def\CH{\mathop{\mathrm{CH}}\nolimits}
\def\proetale{{pro\text{-}\acute{e}tale}}
\def\etale{{\acute{e}tale}}
\def\QCoh{\mathit{QCoh}}
\def\Ker{\mathop{\mathrm{Ker}}}
\def\Im{\mathop{\mathrm{Im}}}
\def\Coker{\mathop{\mathrm{Coker}}}
\def\Coim{\mathop{\mathrm{Coim}}}

% Boxtimes
%
\DeclareMathSymbol{\boxtimes}{\mathbin}{AMSa}{"02}

%
% Macros for moduli stacks/spaces
%
\def\QCohstack{\mathcal{QC}\!\mathit{oh}}
\def\Cohstack{\mathcal{C}\!\mathit{oh}}
\def\Spacesstack{\mathcal{S}\!\mathit{paces}}
\def\Quotfunctor{\mathrm{Quot}}
\def\Hilbfunctor{\mathrm{Hilb}}
\def\Curvesstack{\mathcal{C}\!\mathit{urves}}
\def\Polarizedstack{\mathcal{P}\!\mathit{olarized}}
\def\Complexesstack{\mathcal{C}\!\mathit{omplexes}}
% \Pic is the operator that assigns to X its picard group, usage \Pic(X)
% \Picardstack_{X/B} denotes the Picard stack of X over B
% \Picardfunctor_{X/B} denotes the Picard functor of X over B
\def\Pic{\mathop{\mathrm{Pic}}\nolimits}
\def\Picardstack{\mathcal{P}\!\mathit{ic}}
\def\Picardfunctor{\mathrm{Pic}}
\def\Deformationcategory{\mathcal{D}\!\mathit{ef}}

\setlength{\emergencystretch}{3em}
\begin{document}
\title[Stacks Project, Tag 0DL7]{736 --- Pseudo-coherent derived triples glue across affine pushouts along thickenings}
\maketitle
\noindent Copyright (C) 2005--2025 Johan de Jong. Stacks Project authors. Source: \href{https://stacks.math.columbia.edu/tag/0DL7}{Tag 0DL7}.

\noindent Editorial difficulty note (not author text). Difficulty H3: The proof constructs a homotopy fibre of derived direct images, realizes the gluing isomorphism by strictly compatible finite-free complexes over a henselian pair, and patches their terms over a ring fibre product. This chain-level construction of derived effectivity is specialized research-level algebraic geometry..

\noindent Editorial provenance note (not author text). History: excerpt from revision \texttt{a0444\allowbreak 6e57e\allowbreak c1fbc\allowbreak 252a8\allowbreak 71afc\allowbreak ec775\allowbreak 2fb28\allowbreak 07b14}, spaces-pushouts.tex, lines 1364--1491. Reference presentation was adapted; original-author context and core support blocks were retained or added. The mathematical wording is unchanged. Licensed under GNU FDL 1.2 or later; no invariant sections or cover texts. See ../licenses/COPYING. Context blocks reproduce author definitions and supporting results; they are not additional counted problems.

\begin{definition}
\label{derived-definition-cone}
Let $\mathcal{A}$ be an additive category.
Let $f : K^\bullet \to L^\bullet$ be a morphism of
complexes of $\mathcal{A}$. The {\it cone} of $f$
is the complex $C(f)^\bullet$ given by
$C(f)^n = L^n \oplus K^{n + 1}$ and
differential
$$
d_{C(f)}^n =
\left(
\begin{matrix}
d^n_L & f^{n + 1} \\
0 & -d_K^{n + 1}
\end{matrix}
\right)
$$
It comes equipped with canonical morphisms of complexes
$i : L^\bullet \to C(f)^\bullet$ and $p : C(f)^\bullet \to K^\bullet[1]$
induced by the obvious maps $L^n \to C(f)^n \to K^{n + 1}$.
\end{definition}

\begin{situation}
\label{more-algebra-situation-module-over-fibre-product}
In the following we will consider ring maps
$$
\xymatrix{
B \ar[r] & A & A' \ar[l]
}
$$
where we assume $A' \to A$ is surjective with kernel $I$.
In this situation we set $B' = B \times_A A'$ to
obtain a cartesian square
$$
\xymatrix{
A & A' \ar[l] \\
B \ar[u] & B' \ar[l] \ar[u]
}
$$
\end{situation}

\begin{lemma}
\label{more-algebra-lemma-lift-acyclic-complex}
Let $R$ be a ring. Let $I \subset R$ be an ideal. Let $\mathcal{P}$
be a class of $R$-modules. Assume
\begin{enumerate}
\item each $P \in \mathcal{P}$ is a projective $R$-module,
\item if $P_1 \in \mathcal{P}$ and $P_1 \oplus P_2 \in \mathcal{P}$, then
$P_2 \in \mathcal{P}$, and
\item if $f : P_1 \to P_2$, $P_1, P_2 \in \mathcal{P}$ is surjective
modulo $I$, then $f$ is surjective.
\end{enumerate}
Then given any bounded above acyclic complex $E^\bullet$
whose terms are of the form $P/IP$ for $P \in \mathcal{P}$ there
exists a bounded above acyclic complex $P^\bullet$ whose terms
are in $\mathcal{P}$ lifting $E^\bullet$.
\end{lemma}

\begin{proof}
Say $E^i = 0$ for $i > b$. Assume given $n$ and a morphism of complexes
$$
\xymatrix{
& & P^n \ar[r] \ar[d] & P^{n + 1} \ar[r] \ar[d] & \ldots \ar[r]
& P^b \ar[r] \ar[d] & 0 \ar[r] \ar[d] & \ldots \\
\ldots \ar[r] & E^{n - 1} \ar[r] &
E^n \ar[r] & E^{n + 1} \ar[r] & \ldots \ar[r] &
E^b \ar[r] & 0 \ar[r] & \ldots
}
$$
with $P^i \in \mathcal{P}$, with
$P^n \to P^{n + 1} \to \ldots \to P^b$ acyclic in degrees $\geq n + 1$,
and with vertical maps inducing isomorphisms $P^i/IP^i \to E^i$.
In this situation one can inductively choose isomorphisms
$P^i = Z^i \oplus Z^{i + 1}$ such that the maps $P^i \to P^{i + 1}$
are given by
$Z^i \oplus Z^{i + 1} \to Z^{i + 1} \to Z^{i + 1} \oplus Z^{i + 2}$.
By property (2) and arguing inductively we see that $Z^i \in \mathcal{P}$.
Choose $P^{n - 1} \in \mathcal{P}$ and an isomorphism
$P^{n - 1}/IP^{n - 1} \to E^{n - 1}$. Since
$P^{n - 1}$ is projective and since $Z^n/IZ^n = \Im(E^{n - 1} \to E^n)$,
we can lift the map $P^{n - 1} \to E^{n - 1} \to E^n$ to a map
$P^{n - 1} \to Z^n$. By property (3) the map $P^{n - 1} \to Z^n$
is surjective. Thus we obtain an extension of the diagram
by adding $P^{n - 1}$ and the maps just constructed to the left
of $P^n$. Since a diagram of the desired form exists for $n > b$
we conclude by induction on $n$.
\end{proof}


\begin{lemma}
\label{more-algebra-lemma-lift-complex}
Let $R$ be a ring. Let $I \subset R$ be an ideal. Let $\mathcal{P}$
be a class of $R$-modules.
Let $K \in D(R)$ and let $E^\bullet$ be a complex of $R/I$-modules
representing $K \otimes_R^\mathbf{L} R/I$. Assume
\begin{enumerate}
\item each $P \in \mathcal{P}$ is a projective $R$-module,
\item $P_1 \in \mathcal{P}$ and $P_1 \oplus P_2 \in \mathcal{P}$
if and only if $P_1, P_2 \in \mathcal{P}$,
\item if $f : P_1 \to P_2$, $P_1, P_2 \in \mathcal{P}$ is surjective
modulo $I$, then $f$ is surjective,
\item $E^\bullet$ is bounded above and $E^i$ is of the form $P/IP$
for $P \in \mathcal{P}$, and
\item $K$ can be represented by a bounded above complex
whose terms are in $\mathcal{P}$.
\end{enumerate}
Then there exists a bounded above complex $P^\bullet$ whose terms
are in $\mathcal{P}$ with $P^\bullet/IP^\bullet$ isomorphic to
$E^\bullet$ and representing $K$ in $D(R)$.
\end{lemma}

\begin{proof}
By assumption (5) we can represent $K$ by a bounded above complex 
$K^\bullet$ whose terms are in $\mathcal{P}$. Then $K \otimes_R^\mathbf{L} R/I$
is represented by $K^\bullet/IK^\bullet$. Since $E^\bullet$ is
a bounded above complex of projective $R/I$-modules by (4), 
we can choose a quasi-isomorphism $\delta : E^\bullet \to K^\bullet/IK^\bullet$
(Derived Categories, Lemma
\href{https://stacks.math.columbia.edu/tag/064B}{lemma morphisms from projective complex (Tag 064B)}).
Let $C^\bullet$ be cone on $\delta$
(Derived Categories, Definition \ref{derived-definition-cone}).
The module $C^i$ is the direct sum $K^i/IK^i \oplus E^{i + 1}$
hence is of the form $P/IP$ for some $P \in \mathcal{P}$
as (2) says in particular that $\mathcal{P}$ is preserved under taking sums.
Since $C^\bullet$ is acyclic, we can apply
Lemma \ref{more-algebra-lemma-lift-acyclic-complex} and find a acyclic
lift $A^\bullet$ of $C^\bullet$. The complex $A^\bullet$ is
bounded above and has terms in $\mathcal{P}$. In
$$
\xymatrix{
K^\bullet \ar@{..>}[r] \ar[d] & A^\bullet \ar[d] \\
K^\bullet/IK^\bullet \ar[r] & C^\bullet \ar[r] & E^\bullet[1]
}
$$
we can find the dotted arrow making the diagram commute
by Derived Categories, Lemma
\href{https://stacks.math.columbia.edu/tag/0649}{lemma morphisms lift projective (Tag 0649)}.
We will show below that it follows from (1), (2), (3)
that $K^i \to A^i$ is the inclusion of a direct summand
for every $i$. By property (2) we see that $P^i = \Coker(K^i \to A^i)$
is in $\mathcal{P}$. Thus we can take
$P^\bullet = \Coker(K^\bullet \to A^\bullet)[-1]$ to conclude.

\medskip\noindent
To finish the proof we have to show the following: Let $f : P_1 \to P_2$,
$P_1, P_2 \in \mathcal{P}$ and $P_1/IP_1 \to P_2/IP_2$ is split
injective with cokernel of the form $P_3/IP_3$ for some
$P_3 \in \mathcal{P}$, then $f$ is split injective.
Write $E_i = P_i/IP_i$. Then $E_2 = E_1 \oplus E_3$.
Since $P_2$ is projective we can choose a map $g : P_2 \to P_3$
lifting the map $E_2 \to E_3$. By condition (3) the map $g$
is surjective, hence split as $P_3$ is projective. Set $P_1' = \Ker(g)$
and choose a splitting $P_2 = P'_1 \oplus P_3$. Then $P'_1 \in \mathcal{P}$
by (2). We do not know that
$g \circ f = 0$, but we can consider the map
$$
P_1 \xrightarrow{f} P_2 \xrightarrow{projection} P'_1
$$
The composition modulo $I$ is an isomorphism. Since $P'_1$ is
projective we can split $P_1 = T \oplus P'_1$. If $T = 0$, then
we are done, because then $P_2 \to P'_1$ is a splitting of $f$.
We see that $T \in \mathcal{P}$ by (2).
Calculating modulo $I$ we see that $T/IT = 0$.
Since $0 \in \mathcal{P}$ (as the summand of any $P$ in $\mathcal{P}$)
we see the map $0 \to T$ is surjective and we conclude that $T = 0$
as desired.
\end{proof}


\begin{lemma}
\label{more-algebra-lemma-lift-complex-finite-projectives}
Let $R$ be a ring. Let $I \subset R$ be an ideal. Let $E^\bullet$
be a complex of $R/I$-modules. Let $K$ be an object of $D(R)$. Assume that
\begin{enumerate}
\item $E^\bullet$ is a bounded above complex of finite projective $R/I$-modules,
\item $K \otimes_R^\mathbf{L} R/I$ is represented by $E^\bullet$ in $D(R/I)$,
\item $K$ is pseudo-coherent, and
\item $(R, I)$ is a henselian pair.
\end{enumerate}
Then there exists a bounded above complex $P^\bullet$ of finite projective
$R$-modules representing $K$ in $D(R)$ such that $P^\bullet \otimes_R R/I$
is isomorphic to $E^\bullet$. Moreover, if $E^i$ is free, then $P^i$ is free.
\end{lemma}

\begin{proof}
We apply Lemma \ref{more-algebra-lemma-lift-complex} using the class $\mathcal{P}$
of all finite projective $R$-modules. Properties (1) and (2)
of the lemma are immediate.
Property (3) follows from Nakayama's lemma
(Algebra, Lemma \href{https://stacks.math.columbia.edu/tag/00DV}{lemma NAK (Tag 00DV)}).
Property (4) follows from the fact that we can lift finite projective
$R/I$-modules to finite projective $R$-modules, see
Lemma \ref{more-algebra-lemma-lift-finite-projective-module}.
Property (5) holds because a pseudo-coherent complex can be represented
by a bounded above complex of finite free $R$-modules.
Thus Lemma \ref{more-algebra-lemma-lift-complex} applies and we find $P^\bullet$
as desired. The final assertion of the lemma follows from
Lemma \ref{more-algebra-lemma-isomorphic-finite-projective-lifts}.
\end{proof}


\begin{lemma}
\label{more-algebra-lemma-lift-finite-projective-module}
Let $(R, I)$ be a henselian pair. The map
$$
P \longrightarrow P/IP
$$
induces a bijection between the sets of isomorphism classes of finite
projective $R$-modules and finite projective $R/I$-modules. In particular,
any finite projective $R/I$-module is isomorphic to $P/IP$
for some finite projective $R$-module $P$.
\end{lemma}

\begin{proof}
We first prove the final statement. Let $\overline{P}$ be a finite
projective $R/I$-module. We can find a finite projective module $P'$
over some $R'$ \'etale over $R$ with $R/I = R'/IR'$ such that
$P'/IP'$ is isomorphic to $\overline{P}$, see
Lemma \ref{more-algebra-lemma-lift-projective-module}.
Then, since $(R, I)$ is a henselian pair, the \'etale ring map $R \to R'$ has
a section $\tau : R' \to R$ (Lemma \href{https://stacks.math.columbia.edu/tag/09XI}{lemma characterize henselian pair (Tag 09XI)}).
Setting $P = P' \otimes_{R', \tau} R$ we conclude that $P/IP$
is isomorphic to $\overline{P}$. Of course, this tells us that
the map in the statement of the lemma is surjective.

\medskip\noindent
Injectivity. Suppose that $P_1$ and $P_2$ are finite projective $R$-modules such
that $P_1/IP_1 \cong P_2/IP_2$ as $R/I$-modules. Since $P_1$ is projective,
we can find an $R$-module map $u : P_1 \to P_2$ lifting the given isomorphism.
Then $u$ is surjective by Nakayama's lemma
(Algebra, Lemma \href{https://stacks.math.columbia.edu/tag/00DV}{lemma NAK (Tag 00DV)}).
We similarly find a surjection $v : P_2 \to P_1$.
By Algebra, Lemma \href{https://stacks.math.columbia.edu/tag/05G8}{lemma fun (Tag 05G8)} the map $v \circ u$
is an isomorphism and we conclude $u$ is an isomorphism.
\end{proof}


\begin{lemma}
\label{more-algebra-lemma-lift-projective-module}
Let $A$ be a ring, let $I \subset A$ be an ideal.
Let $\overline{P}$ be a finite projective $A/I$-module.
Then there exists an \'etale ring map $A \to A'$ which induces
an isomorphism $A/I \to A'/IA'$ and a finite projective
$A'$-module $P'$ lifting $\overline{P}$.
\end{lemma}

\begin{proof}
We can choose an integer $n$ and a direct sum decomposition
$(A/I)^{\oplus n} = \overline{P} \oplus \overline{K}$
for some $R/I$-module $\overline{K}$. Choose a lift
$\varphi : A^{\oplus n} \to A^{\oplus n}$ of the projector $\overline{p}$
associated to the direct summand $\overline{P}$.
Let $f \in A[x]$ be the characteristic polynomial of $\varphi$.
Set $B = A[x]/(f)$. By Cayley-Hamilton
(Algebra, Lemma \href{https://stacks.math.columbia.edu/tag/00DX}{lemma charpoly (Tag 00DX)}) there is a map
$B \to \text{End}_A(A^{\oplus n})$ mapping $x$ to $\varphi$.
For every prime $\mathfrak p \supset I$ the image of $f$ in
$\kappa(\mathfrak p)$ is $(x - 1)^rx^{n - r}$ where $r$ is the
dimension of $\overline{P} \otimes_{A/I} \kappa(\mathfrak p)$.
Hence $(x - 1)^nx^n$ maps to zero in $B \otimes_A \kappa(\mathfrak p)$
for all $\mathfrak p \supset I$. Thus $x(1 - x)$ is contained
in every prime ideal of $B/IB$. Hence $x^N(1 - x)^N$ is
contained in $IB$ for some $N \geq 1$.
It follows that $x^N + (1 - x)^N$ is a unit in $B/IB$ and that
$$
\overline{e} = \text{image of }\frac{x^N}{x^N + (1 - x)^N}\text{ in }B/IB
$$
is an idempotent as both assertions hold in $\mathbf{Z}[x]/(x^N(x - 1)^N)$.
The image of $\overline{e}$ in $\text{End}_{A/I}((A/I)^{\oplus n})$ is
$$
\frac{\overline{p}^N}{\overline{p}^N + (1 - \overline{p})^N} = \overline{p}
$$
as $\overline{p}$ is an idempotent. After replacing $A$ by an \'etale
extension $A'$ as in the lemma, we may assume there exists an idempotent
$e \in B$ which maps to $\overline{e}$ in $B/IB$, see
Lemma \ref{more-algebra-lemma-lift-idempotent-upstairs}.
Then the image of $e$ under the map
$$
B = A[x]/(f) \longrightarrow \text{End}_A(A^{\oplus n}).
$$
is an idempotent element $p$ which lifts $\overline{p}$.
Setting $P = \Im(p)$ we win.
\end{proof}


\begin{lemma}
\label{more-algebra-lemma-lift-idempotent-upstairs}
Let $A$ be a ring, let $I \subset A$ be an ideal.
Let $A \to B$ be an integral ring map.
Let $\overline{e} \in B/IB$ be an idempotent.
Then there exists an \'etale ring map $A \to A'$
which induces an isomorphism $A/I \to A'/IA'$ and an idempotent
$e' \in B \otimes_A A'$ lifting $\overline{e}$.
\end{lemma}

\begin{proof}
Choose an element $y \in B$ lifting $\overline{e}$.
Choose $f \in A[x]$ as in Lemma \href{https://stacks.math.columbia.edu/tag/09XG}{lemma helper integral (Tag 09XG)} for $y$.
By Lemma \href{https://stacks.math.columbia.edu/tag/07M1}{lemma lift factorization easy (Tag 07M1)}
we can find an \'etale ring map $A \to A'$ which induces
an isomorphism $A/I \to A'/IA'$ and such that $f = gh$
in $A[x]$ with $g(x) = x^d \bmod IA'$ and $h(x) = (x - 1)^d \bmod IA'$.
After replacing $A$ by $A'$ we may assume that the factorization
is defined over $A$. In that case we see that
$b_1 = g(y) \in B$ is a lift of $\overline{e}^d = \overline{e}$ and
$b_2 = h(y) \in B$ is a lift of
$(\overline{e} - 1)^d = (-1)^d (1 - \overline{e})^d = (-1)^d(1 - \overline{e})$
and moreover $b_1b_2 = 0$. Thus $(b_1, b_2)B/IB = B/IB$ and
$V(b_1, b_2) \subset \Spec(B)$ is disjoint from $V(IB)$. Since
$\Spec(B) \to \Spec(A)$ is closed (see
Algebra, Lemmas \href{https://stacks.math.columbia.edu/tag/00GU}{lemma integral going up (Tag 00GU)} and
\href{https://stacks.math.columbia.edu/tag/00HZ}{lemma going up closed (Tag 00HZ)})
we can find an $a \in A$ which maps to an invertible element
of $A/I$ whose image in $B$ lies in $(b_1, b_2)$, see
Lemma \href{https://stacks.math.columbia.edu/tag/07M3}{lemma separate image closed from closed (Tag 07M3)}.
After replacing $A$ by the localization $A_a$ we get that
$(b_1, b_2) = B$. Then $\Spec(B) = D(b_1) \amalg D(b_2)$;
disjoint union because $b_1b_2 = 0$ and covers
$\Spec(B)$ because $(b_1, b_2) = B$. Let $e \in B$ be the idempotent
corresponding to the open and closed subset $D(b_1)$, see
Algebra, Lemma \href{https://stacks.math.columbia.edu/tag/00EE}{lemma disjoint decomposition (Tag 00EE)}.
Since $b_1$ is a lift of $\overline{e}$ and $b_2$ is a
lift of $\pm (1 - \overline{e})$ we conclude that $e$ is
a lift of $\overline{e}$ by the uniqueness statement in
Algebra, Lemma \href{https://stacks.math.columbia.edu/tag/00EE}{lemma disjoint decomposition (Tag 00EE)}.
\end{proof}


\begin{lemma}
\label{more-algebra-lemma-isomorphic-finite-projective-lifts}
Let $R$ be a ring. Let $I \subset R$ be an ideal. Assume that
every element of $1 + I$ is a unit (in other words $I$ is contained
in the Jacobson radical of $R$). If $P$ and $P'$ are finite
projective $R$-modules, then
\begin{enumerate}
\item if $\varphi : P \to P'$ is an $R$-module map inducing an
isomorphism $\overline{\varphi} : P/IP \to P'/IP'$, then $\varphi$
is an isomorphism,
\item if $P/IP \cong P'/IP'$, then $P \cong P'$.
\end{enumerate}
\end{lemma}

\begin{proof}
Proof of (1). As $P'$ is projective as an $R$-module we may
choose a lift $\psi : P' \to P$ of the map
$P' \to P'/IP' \xrightarrow{\overline{\varphi}^{-1}} P/IP$.
By Nakayama's lemma (Algebra, Lemma \href{https://stacks.math.columbia.edu/tag/00DV}{lemma NAK (Tag 00DV)})
$\psi \circ \varphi$ and $\varphi \circ \psi$ are surjective.
Hence these maps are isomorphisms (Algebra, Lemma \href{https://stacks.math.columbia.edu/tag/05G8}{lemma fun (Tag 05G8)}).
Thus $\varphi$ is an isomorphism.

\medskip\noindent
Proof of (2). Choose an isomorphism $P/IP \cong P'/IP'$.
Since $P$ is projective we can choose a lift $\varphi : P \to P'$ of the map
$P \to P/IP \to P'/IP'$. Then $\varphi$ is an isomorphism by (1).
\end{proof}


\begin{lemma}
\label{more-algebra-lemma-module-over-fibre-product}
In Situation \ref{more-algebra-situation-module-over-fibre-product}
the functor (\href{https://stacks.math.columbia.edu/tag/08KI}{equation functor (Tag 08KI)}) has a right adjoint, namely
the functor
$$
F : (N, M', \varphi) \longmapsto N \times_{\varphi, M} M'
$$
where $M = M'/IM'$. Moreover, the composition of $F$ with
(\href{https://stacks.math.columbia.edu/tag/08KI}{equation functor (Tag 08KI)}) is the identity functor on
$\text{Mod}_B \times_{\text{Mod}_A} \text{Mod}_{A'}$. In other words,
setting $N' = N \times_{\varphi, M} M'$ we have
$N' \otimes_{B'} B = N$ and $N' \otimes_{B'} A' = M'$.
\end{lemma}

\begin{proof}
The adjointness statement follows from the more general
Lemma \href{https://stacks.math.columbia.edu/tag/0D2F}{lemma modules (Tag 0D2F)}.
To prove the final assertion, recall that
$B' = B \times_A A'$ and $N' = N \times_{\varphi, M} M'$ and extend
these equalities to 
$$
\vcenter{
\xymatrix{
A & A' \ar[l] & I \ar[l] \\
B \ar[u] & B' \ar[l] \ar[u] & J \ar[l] \ar[u]
}
}
\quad\text{and}\quad
\vcenter{
\xymatrix{
M & M' \ar[l] & K \ar[l] \\
N \ar[u]_\varphi & N' \ar[l] \ar[u] & L \ar[l] \ar[u]
}
}
$$
where $I, J, K, L$ are the kernels of the horizontal maps of the original
diagrams. We present the proof as a sequence of observations:
\begin{enumerate}
\item $K = IM'$ (see statement lemma),
\item $B' \to B$ is surjective with kernel $J$ and $J \to I$ is bijective,
\item $N' \to N$ is surjective with kernel $L$ and $L \to K$ is bijective,
\item $JN' \subset L$,
\item $\Im(N \to M)$ generates $M$ as an $A$-module
(because $N \otimes_B A = M$),
\item $\Im(N' \to M')$ generates $M'$ as an $A'$-module
(because it holds modulo $K$ and $L$ maps isomorphically to $K$),
\item $JN' = L$ (because $L \cong K = I M'$
is generated by images of elements $x n'$ with $x \in I$ and
$n' \in N'$ by the previous statement),
\item $N' \otimes_{B'} B = N$ (because $N = N'/L$, $B = B'/J$, and
the previous statement),
\item there is a map $\gamma : N' \otimes_{B'} A' \to M'$,
\item $\gamma$ is surjective (see above),
\item the kernel of the composition $N' \otimes_{B'} A' \to M' \to M$ 
is generated by elements $l \otimes 1$ and $n' \otimes x$ with
$l \in K$, $n' \in N'$, $x \in I$ (because $M = N \otimes_B A$ by assumption
and because $N' \to N$ and $A' \to A$ are surjective with kernels
$L$ and $I$),
\item any element of $N' \otimes_{B'} A'$ in the submodule generated
by the elements $l \otimes 1$ and $n' \otimes x$ with
$l \in L$, $n' \in N'$, $x \in I$ can be written as $l \otimes 1$
for some $l \in L$ (because $J$ maps isomorphically to $I$ we see
that $n' \otimes x = n'x \otimes 1$ in $N' \otimes_{B'} A'$;
similarly $x n' \otimes a' = n' \otimes xa' = n'(xa') \otimes 1$
in $N' \otimes_{B'} A'$ when $n' \in N'$, $x \in J$ and $a' \in A'$;
since we have seen that $JN' = L$ this proves the assertion),
\item the kernel of $\gamma$ is zero (because by (10) and (11) any element of
the kernel is of the form $l \otimes 1$ with $l \in L$ which
is mapped to $l \in K \subset M'$ by $\gamma$).
\end{enumerate}
This finishes the proof.
\end{proof}


\begin{lemma}
\label{more-algebra-lemma-flat-module-over-fibre-product}
With $A, A', B, B', I$ as in
Situation \ref{more-algebra-situation-module-over-fibre-product}.
\begin{enumerate}
\item Let $(N, M', \varphi)$ be an object of
$\text{Mod}_B \times_{\text{Mod}_A} \text{Mod}_{A'}$.
If $M'$ is flat over $A'$ and $N$ is flat over $B$, then
$N' = N \times_{\varphi, M} M'$ is flat over $B'$.
\item If $L'$ is a flat $B'$-module, then
$L' = (L \otimes_{B'} B) \times_{(L \otimes_{B'} A)} (L \otimes_{B'} A')$.
\item The category of flat $B'$-modules is equivalent to the
full subcategory of $\text{Mod}_B \times_{\text{Mod}_A} \text{Mod}_{A'}$
consisting of triples
$(N, M', \varphi)$ with $N$ flat over $B$ and $M'$ flat over $A'$.
\end{enumerate}
\end{lemma}

\begin{proof}
In the proof we will use Lemma \ref{more-algebra-lemma-module-over-fibre-product}
without further mention.

\medskip\noindent
Proof of (1). Set $J = \Ker(B' \to B)$. This is an ideal of $B'$
mapping isomorphically to $I = \Ker(A' \to A)$.
Let $\mathfrak b' \subset B'$ be an ideal.
We have to show that $\mathfrak b' \otimes_{B'} N' \to N'$
is injective, see Algebra, Lemma \href{https://stacks.math.columbia.edu/tag/00HD}{lemma flat (Tag 00HD)}.
We know that
$$
\mathfrak b'/(\mathfrak b' \cap J) \otimes_{B'} N' =
\mathfrak b'/(\mathfrak b' \cap J) \otimes_B N \to N
$$
is injective as $N$ is flat over $B$. As
$\mathfrak b' \cap J \to \mathfrak b' \to
\mathfrak b'/(\mathfrak b' \cap J) \to 0$ is exact, we
conclude that it suffices to show that
$(\mathfrak b' \cap J) \otimes_{B'} N' \to N'$
is injective. Thus we may assume that $\mathfrak b' \subset J$.
Next, since $J \to I$ is an isomorphism we have
$$
J \otimes_{B'} N' =
I \otimes_{A'} A' \otimes_{B'} N' =
I \otimes_{A'} M'
$$
which maps injectively into $M'$ as $M'$ is a flat $A'$-module.
Hence $J \otimes_{B'} N' \to N'$ is injective and we conclude
that $\text{Tor}_1^{B'}(B'/J, N') = 0$, see
Algebra, Remark \href{https://stacks.math.columbia.edu/tag/00M6}{remark Tor ring mod ideal (Tag 00M6)}.
Thus we may apply Algebra, Lemma \href{https://stacks.math.columbia.edu/tag/051C}{lemma what does it mean (Tag 051C)}
to $N'$ over $B'$ and the ideal $J$.
Going back to our ideal $\mathfrak b' \subset J$, let
$\mathfrak b' \subset \mathfrak b'' \subset J$ be the smallest ideal
whose image in $I$ is an $A'$-submodule of $I$. In other words,
we have $\mathfrak b'' = A' \mathfrak b'$ if we view $J = I$
as $A'$-module. Then $\mathfrak b''/\mathfrak b'$ is killed
by $J$ and we get a short exact sequence
$$
0 \to \mathfrak b' \otimes_{B'} N' \to
\mathfrak b'' \otimes_{B'} N' \to
\mathfrak b''/\mathfrak b' \otimes_{B'} N' \to 0
$$
by the vanishing of $\text{Tor}_1^{B'}(\mathfrak b''/\mathfrak b', N')$
we get from the application of the lemma.
Thus we may replace $\mathfrak b'$ by $\mathfrak b''$.
In particular we may assume $\mathfrak b'$ is an $A'$-module
and maps to an ideal of $A'$. Then
$$
\mathfrak b' \otimes_{B'} N' =
\mathfrak b' \otimes_{A'} A' \otimes_{B'} N' =
\mathfrak b' \otimes_{A'} M'
$$
This tensor product maps injectively into $M'$ by our assumption that
$M'$ is flat over $A'$. We conclude that
$\mathfrak b' \otimes_{B'} N' \to N' \to M'$ is injective
and hence the first map is injective as desired.

\medskip\noindent
Proof of (2). This follows by tensoring the short exact sequence
$0 \to B' \to B \oplus A' \to A \to 0$ with $L'$ over $B'$.

\medskip\noindent
Proof of (3). Immediate consequence of (1) and (2).
\end{proof}


\begin{lemma}
\label{more-algebra-lemma-finite-module-over-fibre-product}
Let $A, A', B, B', I, M, M', N, \varphi$ be as in
Lemma \ref{more-algebra-lemma-module-over-fibre-product}.
If $N$ finite over $B$ and $M'$ finite over $A'$, then
$N' = N \times_{\varphi, M} M'$ is finite over $B'$.
\end{lemma}

\begin{proof}
We will use the results of
Lemma \ref{more-algebra-lemma-module-over-fibre-product}
without further mention. Choose generators $y_1, \ldots, y_r$ of $N$ over $B$
and generators $x_1, \ldots, x_s$ of $M'$ over $A'$. Using that
$N = N' \otimes_{B'} B$ and $B' \to B$ is surjective we can find
$u_1, \ldots, u_r \in N'$ mapping to $y_1, \ldots, y_r$ in $N$.
Using that $M' = N' \otimes_{B'} A'$ we can find $v_1, \ldots, v_t \in N'$
such that $x_i = \sum v_j \otimes a'_{ij}$ for some $a'_{ij} \in A'$.
In particular we see that the images $\overline{v}_j \in M'$
of the $v_j$ generate $M'$ over $A'$.
We claim that $u_1, \ldots, u_r, v_1, \ldots, v_t$
generate $N'$ as a $B'$-module. Namely, pick $\xi \in N'$. We first choose
$b'_1, \ldots, b'_r \in B'$ such that $\xi$ and $\sum b'_i u_i$ map
to the same element of $N$. This is possible because $B' \to B$
is surjective and $y_1, \ldots, y_r$ generate $N$ over $B$.
The difference $\xi - \sum b'_i u_i$ is of the form $(0, \theta)$
for some $\theta$ in $IM'$. Say $\theta$ is $\sum t_j\overline{v}_j$
with $t_j \in I$. As $J = \Ker(B' \to B)$ maps isomorphically to $I$
we can choose $s_j \in J \subset B'$ mapping to $t_j$.
Because $N' = N \times_{\varphi, M} M'$ it follows
that $\xi = \sum b'_i u_i + \sum s_j v_j$ as desired.
\end{proof}


\begin{lemma}
\label{more-algebra-lemma-finitely-presented-module-over-fibre-product}
Let $A, A', B, B', I$ be as in
Situation \ref{more-algebra-situation-module-over-fibre-product}.
The category of finite projective $B'$-modules
is equivalent to the full subcategory
of $\text{Mod}_B \times_{\text{Mod}_A} \text{Mod}_{A'}$
consisting of triples $(N, M', \varphi)$
with $N$ finite projective over $B$ and $M'$ finite projective over $A'$.
\end{lemma}

\begin{proof}
Recall that a module is finite projective if and only if
it is finitely presented and flat, see
Algebra, Lemma \href{https://stacks.math.columbia.edu/tag/00NX}{lemma finite projective (Tag 00NX)}.
Using Lemmas \ref{more-algebra-lemma-flat-module-over-fibre-product} and
\ref{more-algebra-lemma-finite-module-over-fibre-product}
we reduce to showing that $N' = N \times_{\varphi, M} M'$ is
a $B'$-module of finite presentation if
$N$ finite projective over $B$
and $M'$ finite projective over $A'$.

\medskip\noindent
By Lemma \ref{more-algebra-lemma-finite-module-over-fibre-product}
the module $N'$ is finite over $B'$. Choose a surjection
$(B')^{\oplus n} \to N'$ with kernel $K'$. By base change we obtain maps
$B^{\oplus n} \to N$, $(A')^{\oplus n} \to M'$, and $A^{\oplus n} \to M$
with kernels $K_B$, $K_{A'}$, and $K_A$. There is a canonical map
$$
K' \longrightarrow K_B \times_{K_A} K_{A'}
$$
On the other hand, since $N' = N \times_{\varphi, M} M'$ and
$B' = B \times_A A'$ there is also a
canonical map $K_B \times_{K_A} K_{A'} \to K'$ inverse to the displayed
arrow. Hence the displayed map is an isomorphism. By
Algebra, Lemma \href{https://stacks.math.columbia.edu/tag/0519}{lemma extension (Tag 0519)}
the modules $K_B$ and $K_{A'}$ are finite. We conclude from
Lemma \ref{more-algebra-lemma-finite-module-over-fibre-product}
that $K'$ is a finite $B'$-module provided that $K_B \to K_A$ and
$K_{A'} \to K_A$ induce isomorphisms
$K_B \otimes_B A = K_A = K_{A'} \otimes_{A'} A$.
This is true because the flatness assumptions implies the sequences
$$
0 \to K_B \to B^{\oplus n} \to N \to 0
\quad\text{and}\quad
0 \to K_{A'} \to (A')^{\oplus n} \to M' \to 0
$$
stay exact upon tensoring, see
Algebra, Lemma \href{https://stacks.math.columbia.edu/tag/00HL}{lemma flat tor zero (Tag 00HL)}.
\end{proof}


\begin{lemma}
\label{lemma-pushout-along-thickening-derived}
Let $S$ be a scheme. Consider a pushout
$$
\xymatrix{
X \ar[r]_i \ar[d]_f & X' \ar[d]^{f'}
\\
Y \ar[r]^j & Y'
}
$$
in the category of algebraic spaces over $S$
as in Lemma \href{https://stacks.math.columbia.edu/tag/07VX}{lemma pushout along thickening (Tag 07VX)}.
Assume $i$ is a thickening. Then the essential
image of the functor\footnote{All functors given by derived pullback.}
$$
D(\mathcal{O}_{Y'}) \longrightarrow
D(\mathcal{O}_Y) \times_{D(\mathcal{O}_X)} D(\mathcal{O}_{X'})
$$
contains every triple $(M, K', \alpha)$ where $M \in D(\mathcal{O}_Y)$
and $K' \in D(\mathcal{O}_{X'})$ are pseudo-coherent.
\end{lemma}

\begin{proof}
Let $(M, K', \alpha)$ be an object of the target of the functor
of the lemma. Here $\alpha : Lf^*M \to Li^*K'$
is an isomorphism which is adjoint to a map $\beta : M \to Rf_*Li^*K'$.
Thus we obtain maps
$$
Rj_*M \xrightarrow{Rj_*\beta}
Rj_*Rf_*Li^*K' = Rf'_*Ri_*Li^*K' \leftarrow Rf'_*K'
$$
where the arrow pointing left comes from $K' \to Ri_*Li^*K'$.
Choose a distinguished triangle
$$
M' \to Rj_*M \oplus Rf'_*K' \to Rj_*Rf_*Li^*K' \to M'[1]
$$
in $D(\mathcal{O}_{Y'})$. The first arrow defines canonical maps
$Lj^*M' \to M$ and $L(f')^*M' \to K'$ compatible with $\alpha$.
Thus it suffices to show that the maps
$Lj^*M' \to M$ and $L(f')^*M' \to K$ are isomorphisms.
This we may check \'etale locally on $Y'$, hence we may
assume $Y'$ is \'etale.

\medskip\noindent
Assume $Y'$ affine and $M \in D(\mathcal{O}_Y)$
and $K' \in D(\mathcal{O}_{X'})$ are pseudo-coherent.
Say our pushout corresponds to the fibre product
$$
\xymatrix{
B & B' \ar[l] \\
A \ar[u] & A' \ar[l] \ar[u]
}
$$
of rings where $B' \to B$ is surjective with locally nilpotent kernel $I$
(and hence $A' \to A$ is surjective with locally nilpotent kernel $I$ as well).
The assumption on $M$ and $K'$ imply that $M$ comes from a pseudo-coherent
object of $D(A)$ and $K'$ comes from a pseudo-coherent object of $D(B')$, see
Derived Categories of Spaces, Lemmas
\href{https://stacks.math.columbia.edu/tag/08JL}{lemma pseudo coherent (Tag 08JL)},
\href{https://stacks.math.columbia.edu/tag/071Q}{lemma derived quasi coherent small etale site (Tag 071Q)}, and
\href{https://stacks.math.columbia.edu/tag/08HE}{lemma descend pseudo coherent (Tag 08HE)}
and
Derived Categories of Schemes, Lemma
\href{https://stacks.math.columbia.edu/tag/06Z0}{lemma affine compare bounded (Tag 06Z0)} and
\href{https://stacks.math.columbia.edu/tag/08E7}{lemma pseudo coherent affine (Tag 08E7)}.
Moreover, pushforward and derived pullback agree with the
corresponding operations on derived categories of modules, see
Derived Categories of Spaces, Remark
\href{https://stacks.math.columbia.edu/tag/08GH}{remark match total direct images (Tag 08GH)}
and
Derived Categories of Schemes, Lemmas
\href{https://stacks.math.columbia.edu/tag/0DJK}{lemma quasi coherence pushforward (Tag 0DJK)} and
\href{https://stacks.math.columbia.edu/tag/08DW}{lemma quasi coherence pullback (Tag 08DW)}.
This reduces us to the statement formulated in the next paragraph.
(To be sure these references show
the object $M'$ lies $D_\QCoh(\mathcal{O}_{Y'})$
as this is a triangulated subcategory of $D(\mathcal{O}_{Y'})$.)

\medskip\noindent
Given a diagram of rings as above and a triple
$(M, K', \alpha)$ where $M \in D(A)$, $K' \in D(B')$ are
pseudo-coherent and
$\alpha : M \otimes_A^\mathbf{L} B \to K' \otimes_{B'}^\mathbf{L} B$
is an isomorphism suppose we have distinguished triangle
$$
M' \to M \oplus K' \to K' \otimes_{B'}^\mathbf{L} B \to M'[1]
$$
in $D(A')$. Goal: show that the induced maps
$M' \otimes_{A'}^\mathbf{L} A \to M$ and
$M' \otimes_{A'}^\mathbf{L} B' \to K'$ are isomorphisms.
To do this, choose a bounded above complex
$E^\bullet$ of finite free $A$-modules representing $M$.
Since $(B', I)$ is a henselian pair
(More on Algebra, Lemma \href{https://stacks.math.columbia.edu/tag/0ALI}{lemma locally nilpotent henselian (Tag 0ALI)})
with $B = B'/I$ we may apply More on Algebra, Lemma
\ref{more-algebra-lemma-lift-complex-finite-projectives}
to see that there exists a bounded above complex $P^\bullet$
of free $B'$-modules such that $\alpha$ is represented
by an isomorphism $E^\bullet \otimes_A B \cong P^\bullet \otimes_{B'} B$.
Then we can consider the short exact sequence
$$
0 \to L^\bullet \to
E^\bullet \oplus P^\bullet \to P^\bullet \otimes_{B'} B \to 0
$$
of complexes of $B'$-modules.
More on Algebra, Lemma
\ref{more-algebra-lemma-finitely-presented-module-over-fibre-product}
implies $L^\bullet$ is a bounded above complex of
finite projective $A'$-modules
(in fact it is rather easy to show directly that $L^n$ is finite free
in our case) and that we have
$L^\bullet \otimes_{A'} A = E^\bullet$ and
$L^\bullet \otimes_{A'} B' = P^\bullet$.
The short exact sequence gives a distinguished triangle
$$
L^\bullet \to M \oplus K' \to K' \otimes_{B'}^\mathbf{L} B \to (L^\bullet)[1]
$$
in $D(A')$ (Derived Categories, Section
\href{https://stacks.math.columbia.edu/tag/014Z}{section canonical delta functor (Tag 014Z)}) which is isomorphic
to the given distinguished triangle by general properties of
triangulated categories (Derived Categories, Section
\href{https://stacks.math.columbia.edu/tag/05QN}{section elementary results (Tag 05QN)}). In other words, $L^\bullet$
represents $M'$ compatibly with the given maps. Thus the maps
$M' \otimes_{A'}^\mathbf{L} A \to M$ and
$M' \otimes_{A'}^\mathbf{L} B' \to K'$ are
isomorphisms because we just saw that the corresponding
thing is true for $L^\bullet$.
\end{proof}
\end{document}
